\section{Methods}
\label{subsec:methods}

\subsection{Overview}
Figure~\ref{fig:pipeline_overview} provides an overview of our full pipeline.

\begin{figure*}[t]
    \centering
    \includegraphics[width=\textwidth]{figs/fig_pipeline/pipeline_overview.png}
    \caption{
    Pipeline overview of the proposed method. We first build a stable geometric scaffold from fixed-pose COLMAP sparse support with weak geometric priors, then perform no-densify reliability-aware geometry refinement, and finally freeze the main geometry for YCbCr-decoupled appearance refinement with calibrated proxy supervision.}
    \label{fig:pipeline_overview}
\end{figure*}

\begin{table*}[t]
\setlength{\tabcolsep}{3pt}
\renewcommand{\arraystretch}{1.08}
\footnotesize
\centering
\begin{tabular}{@{}>{\raggedright\arraybackslash}p{0.18\textwidth}c>{\raggedright\arraybackslash}p{0.22\textwidth}>{\raggedright\arraybackslash}p{0.18\textwidth}>{\raggedright\arraybackslash}p{0.23\textwidth}@{}}
\toprule
Cue / Supervision & Stage & Why Reliable & How Used & Role \\
\midrule
Weak monocular depth & S1--S2 & coarse depth trend is more reliable than absolute scale & low weight + patch sampling + annealing & reduce large-scale geometric ambiguity without over-constraining the scene \\
Rendered multi-view depth consistency & S1 & cross-view agreement indicates stable geometry during densification & valid-overlap masking + geometric-depth thresholding & suppress unstable densification and early floating artifacts \\
Illumination-invariant structure prior & S2 & structural edges remain more stable than raw RGB under low light & confidence-modulated structural supervision & preserve contours and layout when photometric supervision is unreliable \\
Sparse-guided COLMAP support & S2 & longer tracks, lower reprojection error, and denser neighborhoods are more trustworthy & quality/density-aware barycentric support weighting & prevent geometry drift during no-densify refinement \\
Luminance proxy supervision & S3 & luminance recovery is more reliable with a shadow-aware calibrated proxy than direct RGB regression & proxy target + $W_Y$ reweighting & recover exposure conservatively, especially in severely dark regions \\
Chroma correction supervision & S3 & chroma is relatively more reliable outside severe shadows but still requires conservative correction & supervision-image chroma + $W_C$ reweighting + residual regularization & correct color cast without introducing excessive chromatic shift \\
\bottomrule
\end{tabular}
\caption{Reliability-aware supervision allocation across stages.}
\label{tab:reliability_summary}
\end{table*}

Our method follows a three-stage pipeline: geometry bootstrap, reliability-aware
geometry refinement, and geometry-frozen appearance refinement. In the first
stage, we build a stable geometric scaffold from fixed-pose COLMAP sparse
points after filtering and voxel deduplication, optionally together with a
small number of random Gaussians for exploration. At this stage, monocular
depth is only used as a weak geometric trend prior, while a lightweight
multi-view reprojection consistency is introduced to regularize rendered
geometric depth during densification. The first-stage objective can be summarized as
\begin{equation}
\mathcal{L}^{(1)}
= \mathcal{L}_{\mathrm{rgb}}
+ \lambda_d^{(1)} \mathcal{L}_{\mathrm{depth}}
+ \lambda_{mv} \mathcal{L}_{\mathrm{mv}}
+ \lambda_{exp} \mathcal{L}_{\mathrm{exp}},
\end{equation}
where $\mathcal{L}_{\mathrm{rgb}}$ aligns the base rendering with the enhanced
supervision image, $\mathcal{L}_{\mathrm{depth}}$ provides weak depth guidance,
and $\mathcal{L}_{\mathrm{mv}}$ enforces rendered geometric depth consistency
across neighboring views. This stage is also the only stage that performs densification. 

In the second stage, we warm-start from the first-stage geometry and continue
no-densify refinement with weak depth, an illumination-invariant structure prior,
and a quality- and density-aware sparse-guided regularization that reuses COLMAP
sparse geometry as persistent training-time support. The multi-view reprojection
term is removed, and the weak depth prior is gradually annealed as geometry becomes
more stable. The Stage 2 objective is
\begin{equation}
\mathcal{L}^{(2)}
= \mathcal{L}_{\mathrm{rgb}}
+ \lambda_d^{(2)}(t)\,\mathcal{L}_{\mathrm{depth}}
+ \lambda_s \mathcal{L}_{\mathrm{struct}}
+ \lambda_{sp} \mathcal{L}_{\mathrm{sparse}},
\end{equation}
where $\lambda_d^{(2)}(t)$ decreases over training, and $\mathcal{L}_{\mathrm{sparse}}$
uses sparse-point quality and local density to form reliability-aware soft geometric support. 

In the final stage, we freeze the main geometric parameters and perform geometry-frozen
refinement in YCbCr space. Specifically, the illumination branch only restores
luminance in the Y channel under a calibrated proxy target, while an additive chroma
residual branch corrects Cb/Cr against the more reliable supervision image. The Y and
Cb/Cr reconstruction terms are further modulated by independent weight maps, yielding
a reliability-aware allocation of supervision across both stages and image regions.
The Stage 3 objective is written as
\begin{equation}
\mathcal{L}^{(3)}
= \mathcal{L}_{\mathrm{rgb}}
+ \lambda_Y \mathcal{L}_{Y}
+ \lambda_{CbCr} \mathcal{L}_{CbCr}
+ \lambda_{chr} \|\Delta_C\|_1,
\end{equation}
where $\mathcal{L}_{Y}$ aligns $Y(\hat{I}_{rec})$ with the calibrated proxy target,
$\mathcal{L}_{CbCr}$ aligns $CbCr(\hat{I}_{rec})$ with the supervision image
under an independent chroma weight map, and $\|\Delta_C\|_1$
regularizes the additive chroma residual. 

\subsection{Sparse-Guided Regularization}
To stabilize geometry refinement in Stage 2, we reuse the sparse points from Structure-from-Motion(SfM) as persistent training-time support. For each sampled active Gaussian center $\mu_i$, we retrieve its $K$ nearest sparse points and regularize $\mu_i$ toward a local support-weighted barycenter, rather than directly snapping it to a single nearest point. This design reduces the sensitivity to isolated sparse outliers and makes the constraint depend on locally consistent sparse support.

We define the support score $s_j$ of each sparse point $p_j$ as
\begin{equation}
s_j = q_j \cdot d_j,
\end{equation}
where $q_j$ denotes the observation-quality term and $d_j$ denotes the local-density term.

The observation-quality term is defined as
\begin{equation}
q_j = \phi(t_j)\,\psi(e_j),
\end{equation}
where $t_j$ and $e_j$ are respectively the track length and reprojection error associated with sparse point $p_j$. In our implementation,
\begin{equation}
\phi(t_j)=\log(1+t_j), \qquad
\psi(e_j)=\exp\!\left(-\frac{e_j}{\eta}\right),
\end{equation}
with $\eta$ set to the median reprojection error. As a result, sparse points supported by longer tracks and smaller reprojection errors receive larger quality scores.

To further favor sparse points lying in locally well-supported regions, we define a density term
\begin{equation}
d_j = \frac{\bar{r}}{r_j+\varepsilon},
\end{equation}
where $r_j$ is the average distance from $p_j$ to its local sparse neighbors, and $\bar{r}$ is the median of these local radii over the sparse cloud. This term increases the contribution of sparse points located in relatively denser and more stable local neighborhoods.

Using $s_j$, the barycentric weight of neighbor $p_j$ for Gaussian $\mu_i$ is
\begin{equation}
w_{ij}=
\frac{\dfrac{s_j}{\|\mu_i-p_j\|_2+\varepsilon}}
{\sum\limits_{l\in \mathcal{N}_K(i)}
\dfrac{s_l}{\|\mu_i-p_l\|_2+\varepsilon}},
\end{equation}
and the sparse support barycenter is
\begin{equation}
\bar{p}_i = \sum_{j\in\mathcal{N}_K(i)} w_{ij} p_j .
\end{equation}
Therefore, each Gaussian is attracted not to an individual sparse point, but to a reliability-aware local support center determined jointly by sparse-point quality, local density, and geometric proximity.

The sparse-guided loss is then
\begin{equation}
\mathcal{L}_{\mathrm{sparse}}
=
\frac{1}{|\mathcal{S}|}
\sum_{i\in\mathcal{S}}
\rho\!\left(\|\mu_i-\bar{p}_i\|_2\right),
\end{equation}
where $\mathcal{S}$ is the sampled set of active Gaussians and $\rho(\cdot)$ is a Charbonnier-style robust penalty:
\begin{equation}
\rho(x)=\sqrt{x^2+c^2}-c.
\end{equation}
This robust form prevents distant mismatches from dominating the optimization, so the sparse guidance remains weak but persistent throughout no-densify geometry refinement.

\begin{figure*}[!t]
\centering
\setlength{\tabcolsep}{0pt}
\newcommand{\FigThreeSceneW}{0.018\textwidth}
\newcommand{\FigThreeColW}{0.1615\textwidth}
\newcommand{\FigThreeGap}{0pt}
\newcommand{\FigThreeMainH}{1.70cm}
\newcommand{\FigThreeSceneLabel}[1]{%
    \rotatebox{90}{\small\bfseries #1}%
}
\newcommand{\FigThreePlain}[1]{%
    \includegraphics[height=\FigThreeMainH]{#1}%
}
\newcommand{\FigThreeColLabel}[1]{%
    \parbox[t][2.1em][t]{\linewidth}{\centering\footnotesize\bfseries #1}%
}

% ----------------------------- Chocolate row -----------------------------
\begin{minipage}[c]{\FigThreeSceneW}
    \centering
    \FigThreeSceneLabel{Chocolate}
\end{minipage}\hspace{\FigThreeGap}%
\begin{minipage}[c]{\FigThreeColW}
    \centering
    \FigThreePlain{figs/fig3/chocolate_lowlight.JPG}
\end{minipage}\hspace{\FigThreeGap}%
\begin{minipage}[c]{\FigThreeColW}
    \centering
    \FigThreePlain{figs/fig3/chocolate_alethnerf.JPG}
\end{minipage}\hspace{\FigThreeGap}%
\begin{minipage}[c]{\FigThreeColW}
    \centering
    \FigThreePlain{figs/fig3/chocolate_i2nerf.JPG}
\end{minipage}\hspace{\FigThreeGap}%
\begin{minipage}[c]{\FigThreeColW}
    \centering
    \FigThreePlain{figs/fig3/chocolate_litags.JPG}
\end{minipage}\hspace{\FigThreeGap}%
\begin{minipage}[c]{\FigThreeColW}
    \centering
    \FigThreePlain{figs/fig3/chocolate_luminancegs.JPG}
\end{minipage}\hspace{\FigThreeGap}%
\begin{minipage}[c]{\FigThreeColW}
    \centering
    \FigThreePlain{figs/fig3/chocolate_ours.JPG}
\end{minipage}

\vspace{0.03em}

% ----------------------------- GearWorks row -----------------------------
\begin{minipage}[c]{\FigThreeSceneW}
    \centering
    \FigThreeSceneLabel{GearWorks}
\end{minipage}\hspace{\FigThreeGap}%
\begin{minipage}[c]{\FigThreeColW}
    \centering
    \FigThreePlain{figs/fig3/gearworks_lowlight.JPG}
\end{minipage}\hspace{\FigThreeGap}%
\begin{minipage}[c]{\FigThreeColW}
    \centering
    \FigThreePlain{figs/fig3/gearworks_alethnerf.JPG}
\end{minipage}\hspace{\FigThreeGap}%
\begin{minipage}[c]{\FigThreeColW}
    \centering
    \FigThreePlain{figs/fig3/gearworks_i2nerf.JPG}
\end{minipage}\hspace{\FigThreeGap}%
\begin{minipage}[c]{\FigThreeColW}
    \centering
    \FigThreePlain{figs/fig3/gearworks_litags.JPG}
\end{minipage}\hspace{\FigThreeGap}%
\begin{minipage}[c]{\FigThreeColW}
    \centering
    \FigThreePlain{figs/fig3/gearworks.luminancegs.JPG}
\end{minipage}\hspace{\FigThreeGap}%
\begin{minipage}[c]{\FigThreeColW}
    \centering
    \FigThreePlain{figs/fig3/gearworks_ours.JPG}
\end{minipage}

\vspace{0.03em}

% ----------------------------- Popcorn row -----------------------------
\begin{minipage}[c]{\FigThreeSceneW}
    \centering
    \FigThreeSceneLabel{Popcorn}
\end{minipage}\hspace{\FigThreeGap}%
\begin{minipage}[c]{\FigThreeColW}
    \centering
    \FigThreePlain{figs/fig3/popcorn_lowlight.JPG}
\end{minipage}\hspace{\FigThreeGap}%
\begin{minipage}[c]{\FigThreeColW}
    \centering
    \FigThreePlain{figs/fig3/popcorn_alethnerf.JPG}
\end{minipage}\hspace{\FigThreeGap}%
\begin{minipage}[c]{\FigThreeColW}
    \centering
    \FigThreePlain{figs/fig3/popcorn_i2nerf.JPG}
\end{minipage}\hspace{\FigThreeGap}%
\begin{minipage}[c]{\FigThreeColW}
    \centering
    \FigThreePlain{figs/fig3/popcorn_litags.JPG}
\end{minipage}\hspace{\FigThreeGap}%
\begin{minipage}[c]{\FigThreeColW}
    \centering
    \FigThreePlain{figs/fig3/popcorn_luminancegs.JPG}
\end{minipage}\hspace{\FigThreeGap}%
\begin{minipage}[c]{\FigThreeColW}
    \centering
    \FigThreePlain{figs/fig3/popcorn_ours.JPG}
\end{minipage}

\vspace{0.03em}

% ----------------------------- Sculpture row -----------------------------
\begin{minipage}[c]{\FigThreeSceneW}
    \centering
    \FigThreeSceneLabel{Sculpture}
\end{minipage}\hspace{\FigThreeGap}%
\begin{minipage}[c]{\FigThreeColW}
    \centering
    \FigThreePlain{figs/fig3/sculpture_lowlight.JPG}
\end{minipage}\hspace{\FigThreeGap}%
\begin{minipage}[c]{\FigThreeColW}
    \centering
    \FigThreePlain{figs/fig3/sculpture_alethnerf.JPG}
\end{minipage}\hspace{\FigThreeGap}%
\begin{minipage}[c]{\FigThreeColW}
    \centering
    \FigThreePlain{figs/fig3/sculpture_i2nerf.JPG}
\end{minipage}\hspace{\FigThreeGap}%
\begin{minipage}[c]{\FigThreeColW}
    \centering
    \FigThreePlain{figs/fig3/sculpture_litags.JPG}
\end{minipage}\hspace{\FigThreeGap}%
\begin{minipage}[c]{\FigThreeColW}
    \centering
    \FigThreePlain{figs/fig3/sculpture_luminancegs.JPG}
\end{minipage}\hspace{\FigThreeGap}%
\begin{minipage}[c]{\FigThreeColW}
    \centering
    \FigThreePlain{figs/fig3/sculpture_ours.JPG}
\end{minipage}

\vspace{0.03em}

% ----------------------------- Ujikintoki row -----------------------------
\begin{minipage}[c]{\FigThreeSceneW}
    \centering
    \FigThreeSceneLabel{Ujikintoki}
\end{minipage}\hspace{\FigThreeGap}%
\begin{minipage}[c]{\FigThreeColW}
    \centering
    \FigThreePlain{figs/fig3/ujikintoki_lowlight.JPG}
\end{minipage}\hspace{\FigThreeGap}%
\begin{minipage}[c]{\FigThreeColW}
    \centering
    \FigThreePlain{figs/fig3/ujikintoki_alethnerf.JPG}
\end{minipage}\hspace{\FigThreeGap}%
\begin{minipage}[c]{\FigThreeColW}
    \centering
    \FigThreePlain{figs/fig3/ujikintoki_i2nerf.JPG}
\end{minipage}\hspace{\FigThreeGap}%
\begin{minipage}[c]{\FigThreeColW}
    \centering
    \FigThreePlain{figs/fig3/ujikintoki_litags.JPG}
\end{minipage}\hspace{\FigThreeGap}%
\begin{minipage}[c]{\FigThreeColW}
    \centering
    \FigThreePlain{figs/fig3/ujikintoki_luminancegs.JPG}
\end{minipage}\hspace{\FigThreeGap}%
\begin{minipage}[c]{\FigThreeColW}
    \centering
    \FigThreePlain{figs/fig3/ujikintoki_ours.JPG}
\end{minipage}

\vspace{0.10em}

\begin{minipage}[c]{\FigThreeSceneW}
    \centering
\end{minipage}\hspace{\FigThreeGap}%
\begin{minipage}[c]{\FigThreeColW}
    \centering
    \FigThreeColLabel{Low-light input}
\end{minipage}\hspace{\FigThreeGap}%
\begin{minipage}[c]{\FigThreeColW}
    \centering
    \FigThreeColLabel{AlethNeRF}
\end{minipage}\hspace{\FigThreeGap}%
\begin{minipage}[c]{\FigThreeColW}
    \centering
    \FigThreeColLabel{I2NeRF}
\end{minipage}\hspace{\FigThreeGap}%
\begin{minipage}[c]{\FigThreeColW}
    \centering
    \FigThreeColLabel{LITA-GS}
\end{minipage}\hspace{\FigThreeGap}%
\begin{minipage}[c]{\FigThreeColW}
    \centering
    \FigThreeColLabel{Luminance-GS}
\end{minipage}\hspace{\FigThreeGap}%
\begin{minipage}[c]{\FigThreeColW}
    \centering
    \FigThreeColLabel{Ours}
\end{minipage}

\vspace{-0.55em}
\caption{Qualitative comparison on five representative low-light scenes.
Each row shows the low-light input, four strong baselines, and our result.
Our method consistently recovers clearer scene structures, more plausible
brightness, and more natural colors under severe low-light conditions.}
\label{fig:qualitative_results}
\end{figure*}


\subsection{Geometry-Frozen Decoupled Appearance Refinement}
After the geometry becomes sufficiently stable, further joint optimization of
geometry and appearance may allow photometric compensation to corrupt the
learned structure. We therefore freeze the main geometric parameters $\Theta_{\mathrm{geom}}$ in the
final stage and restrict optimization to appearance only. Specifically, we
freeze the Gaussian center $\mu$, rotation $r$, and scale $s$,
while continuing to optimize the appearance-related parameters
\begin{equation}
\Theta_{\mathrm{app}}=\{o,c,\phi^{illum},\phi^{chr}\},
\end{equation}
where $o$, $c$, $\phi^{illum}$ and $\phi^{chr}$ denote the opacity, SH color, the illumination head and the chroma head respectively. Given camera $C$, the base
rendering and two auxiliary appearance heads are
\begin{equation}
\begin{gathered}
I_{base}
=
\mathcal{R}_{rgb}
\!\left(
\Theta;C
\right), \\
A_Y
=
\mathcal{R}_{illum}
\!\left(
\Theta_{\mathrm{geom}},\phi^{illum};C
\right), \\
A_C
=
\mathcal{R}_{chr}
\!\left(
\Theta_{\mathrm{geom}},\phi^{chr};C
\right).
\end{gathered}
\end{equation}
This stage is designed around a simple observation: under low-light inputs,
luminance restoration and chromatic correction do not exhibit the same
reliability. We therefore refine appearance in YCbCr space and supervise the
luminance and chroma branches with different targets.

\paragraph{Calibrated proxy target for luminance refinement.}
For a low-light input image, we construct two complementary targets:
an enhanced supervision image $I_{sup}$ and a calibrated proxy target
$I_{proxy}$. The enhanced image serves as a globally plausible appearance
reference, while the proxy target is used only for luminance restoration. To
make the proxy target conservative in already bright regions and more
supportive in severely under-exposed areas, we combine a global exposure
compensation branch $I_{global}$ and a shadow-aware branch $I_{shadow}$ through
a non-linear, smoothed shadow confidence map $w_{sh}$ that approaches 1 in dark areas and 0 in bright areas:
\begin{equation}
\begin{gathered}
I_{proxy}(x)
=
\big(1-w_{sh}(x)\big)\,I_{global}(x) \\
\quad + w_{sh}(x)\,I_{shadow}(x).
\end{gathered}
\end{equation}

This proxy construction allows the luminance branch to follow a
statistically-calibrated and shadow-aware target, instead of directly
overfitting to the enhanced supervision image everywhere.

\paragraph{YCbCr-decoupled appearance modeling.}
We explicitly decouple luminance restoration from chromatic correction in YCbCr
space. Instead of directly relighting all RGB channels together, which may
entangle exposure recovery with unintended color shifts, the illumination branch
only modulates the luminance channel:
\begin{equation}
Y' = \operatorname{clip}(2\,\sigma(A_Y)\odot Y,0,1).
\end{equation}
The relit result is then reconstructed with the updated luminance $Y'$ while
preserving the original chroma channels $(C_b,C_r)$. In this way, brightness
recovery is handled independently of color correction.

On top of the relit result, we apply a separate chroma branch that predicts a
bounded additive residual only on the chroma channels:
\begin{equation}
\begin{gathered}
\Delta_C = \alpha_c\tanh(A_C), \\
\hat C_b = C_b+\Delta_{cb}, \qquad \hat C_r = C_r+\Delta_{cr},
\end{gathered}
\end{equation}
where $\Delta_C=(\Delta_{cb},\Delta_{cr})$ and $\alpha_c$ controls the chroma
residual range. This residual design keeps chromatic correction conservative:
the illumination branch addresses under-exposure, while the chroma branch only
removes remaining color bias without reintroducing large chromatic changes.

\paragraph{Branch-specific supervision and final objective.}
The decoupled refinement is trained with branch-specific supervision. We retain
a base RGB reconstruction loss against the enhanced supervision image $I_{sup}$
so that the underlying 3DGS appearance remains globally plausible. Beyond target assignment mentioned above, we further modulate the supervision spatially using
branch-specific confidence maps. Let $W_Y$ and $W_C$ denote the luminance and
chroma confidence maps, respectively. $W_Y$ emphasizes under-exposed regions
where brightness recovery remains meaningful, whereas $W_C$ down-weights severe
shadow regions where chroma observations are less reliable. In this way, target
assignment determines \emph{which} signal supervises each branch, while
$W_Y$ and $W_C$ determine \emph{where} that supervision should be trusted more.

For luminance restoration, let
$Y_{rec}=Y(I_{rec})$ and $Y_{proxy}=Y(I_{proxy})$. We define
\begin{equation}
\mathcal{L}_{Y}
=
\operatorname{WeightedL1}
\big(
Y_{rec},Y_{proxy};W_Y
\big).
\end{equation}

For chromatic correction, we use a color-cast-aware chroma objective. Since
relighting may still leave a residual low-frequency chroma drift, often
manifesting as a green-dominant cast, we do not
directly regress to the raw supervision chroma. Instead, we define
$(Y^{sup},C_b^{sup},C_r^{sup})=\operatorname{YCbCr}(I_{sup})$ and apply a small
bias to obtain $C_b^{*}=C_b^{sup}+b_{cb}$ and $C_r^{*}=C_r^{sup}+b_{cr}$. We
further use asymmetric channel weights to allow different supervision strengths
for the Cb and Cr components. Let
$\tilde{C}_{rec}=(\omega_{cb}\hat C_b,\omega_{cr}\hat C_r)$ and
$\tilde{C}_{tgt}=(\omega_{cb}C_b^{*},\omega_{cr}C_r^{*})$. The chroma loss is
\begin{equation}
\mathcal{L}_{CbCr}
=
\operatorname{WeightedL1}
\big(
\tilde{C}_{rec},\tilde{C}_{tgt};W_C
\big)
+\lambda_{gm}\mathcal{L}_{CbCr}^{mean}.
\end{equation}

The final-stage objective is
\begin{equation}
\mathcal{L}^{(3)}
=
\mathcal{L}_{rgb}
+\lambda_Y \mathcal{L}_{Y}
+\lambda_{CbCr} \mathcal{L}_{CbCr}
+\lambda_{chr}\|\Delta_C\|_1.
\end{equation}
